\documentclass[12pt, a4paper]{article}
\usepackage[utf8]{inputenc}
\usepackage[italian]{babel}
\usepackage{fancyhdr} %Header e Footer
\usepackage{titling} %Spaziatura Titolo
\usepackage{enumitem}
\usepackage{setspace} 
\usepackage{siunitx} %Notazione scientifica e unità di misura
\usepackage{parskip} %Spaziatura paragrafi automatica
\usepackage{graphicx} %Imaggini
\usepackage{float} 
\usepackage{booktabs}
\usepackage[
    top=3cm,    
    bottom=2cm, 
    left=1.5cm,   
    right=1.5cm,
    headheight=15pt]{geometry} %Margini
\usepackage[colorlinks=true, linkcolor=cyan]{hyperref} %Colore link
\usepackage{amsmath}
\usepackage[italian, capitalize]{cleveref} %Nuovo nome refernce
\usepackage{caption}

\setstretch{0.85}
\title{5° Esperienza Laboratorio di Fisica 2} %NUMERO ESPERIENZA   
\author{Gruppo B02: Bifulco Flavio DC6000147 - De Marco Alessandro DC6000138 - \\ De Simone Adriano DC6000154}
\date{25/11/2026} %DATA
\setlength{\droptitle}{-2cm}

% Titolo in grassetto e riduzione spazio sotto 
\pretitle{\begin{center}\huge\bfseries} 
\posttitle{\par\end{center}} 

% Autore più piccolo 
\preauthor{\begin{center}\small} 
\postauthor{\par\end{center}}

\postdate{
    \par\end{center} 
    \vspace{0.5em}   
    \begin{center}
        \textit{\textbf{\Large Studio in regime impulsivo di circuiti RC e RLC}} %TITOLO
    \end{center}
    \vspace{0.5cm} 
}

\newcommand{\E}{\`E } % Comando per "è" grande

\graphicspath{{./immagini/}}
\renewcommand{\thefigure}{\arabic{section}.\arabic{figure}}
\addto\captionsitalian{\renewcommand{\figurename}{Immagine}}
\crefname{figure}{Immagine}{immagini}
\makeatletter
\@addtoreset{figure}{section}
\makeatother

\pagestyle{fancy}
\fancyhf{}

% Configurazione Header
\fancyhead[L]{Gruppo B02: Bifulco Flavio - De Marco Alessandro - De Simone Adriano}       
\fancyhead[R]{\thedate}                
\renewcommand{\headrulewidth}{0pt}
\addtolength{\headsep}{10pt}

% Configurazione Footer
\fancyfoot[C]{\thepage} 

\sisetup{
    output-decimal-marker = {,},
    separate-uncertainty = true,
    per-mode = fraction
}
\DeclareSIUnit\cal{cal}

\renewcommand{\thesection}{\arabic{section}}
\renewcommand{\thesubsection}{\arabic{section}.\arabic{subsection}}
\renewcommand{\thesubsubsection}{\arabic{subsubsection}}
\renewcommand{\theHsection}{\arabic{section}}
\renewcommand{\theHsubsection}{\arabic{section}.\arabic{subsection}}
\renewcommand{\theHsubsubsection}{\arabic{section}.\arabic{subsection}.\arabic{subsubsection}}
\hypersetup{bookmarksnumbered}
\begin{document}
\maketitle
%\textbf{\Large Obiettivo:}
%\begin{enumerate}[label=\arabic*), leftmargin=2cm]
%    \item Verifica sperimentalmente la prima legge di Ohm con il metodo volt-amperometrico. 
%\end{enumerate}
%
%\textbf{\Large Strumentazione:}
%\begin{itemize}[label= -, leftmargin=2cm]
%    \item Bilancia (\(\Delta G_\text{sens} = \qty{0.1}{\gram}\)).
%    \item 
%\end{itemize}
%
%\bigskip \bigskip
%\noindent{\LARGE \textbf{Procedimento di misura e analisi dati: }}\(\)
%\bigskip

\section{Introduzione}
L'obiettivo dell'esperienza è quello di andare a studiare il comportamento di questi due circuiti notevoli nel 
momento in cui ricevono in ingresso un voltaggio impulsivo. Simuleremo allora attraverso un onda quadra un impulso periodico.
\section{Apparato sperimentale}
\begin{itemize}[leftmargin=2cm]
    \item Generatore arbitrario di funzioni.
    \item Oscilloscopio Teledyne Lecroy HDO 6104B [CONTROLLARE]
    \item RCL meter.
    \item Due blocchi di impedenze per il circuito RC, con valori misurati:
    \begin{itemize}[leftmargin=1cm]
        \item Resistenza da \(R = \qty{100.3(0.35)}{\ohm} \) [MODIFICARE]
        \item Condensatore da \(C = \qty{219.0(0.6)}{\nano\farad} \) [MODIFICARE]
    \end{itemize}
    \item Tre blocchi di impedenze per il circuito RCL, con valori misurati:
    \begin{itemize}[leftmargin=1cm]
        \item Resistenza da \(R = \qty{100.3(0.35)}{\ohm} \) [MODIFICARE]
        \item Induttanza da \(L = \qty{9.891(0.025)}{\milli\henry} \) [MODIFICARE]
        \item Condensatore da \(C = \qty{219.0(0.6)}{\nano\farad} \) [MODIFICARE]
    \end{itemize}
    \item Basetta e cavi:
    \begin{itemize}[leftmargin=1cm]
        \item Cavo BNC-BNC.
        \item Cavi BNC-Coccodrilli.
    \end{itemize}
\end{itemize}


\section{Circuito RC}
\subsection{Introduzione}
Anche in questo caso la risoluzione analitica dell'equazione differenziale del circuito, con le adeguate condizioni iniziali, permette 
di descriverlo completamente.

Nel nostro caso otterremo una tensione ai capi del condensatore dipendente dalla funzione:

 
\[V_c(t)=
\begin{cases}
-V_0+[(V_0+V^*)e^{-\frac{t}{\tau}}] \quad se \: 0\leq t<\frac{T}{2}\\
+V_0-[(V_0+V^*)e^{-\frac{t-\frac{T}{2}}{\tau}}] \quad se \; \frac{T}{2}\leq t<T
\end{cases}
\]

Dove: \[ V^*= V_0 \frac{1-e^{-\frac{T}{2\tau}}}{1+e^{-\frac{T}{2\tau}}}\]
\[\text{e} \quad T \quad\text{è il periodo dell'onda qadra}\]

Cio ci permette allora di fare delle considerazioni, a priori, di come si dovrebbe comportare il circuito reale
\begin{itemize}[leftmargin=2cm]
    \item Per un tempo caratteristico sufficentemente minore del periodo dell'onda quadra, ci aspettiamo che il condensatore riesca a caricarsi completamente
    \item Se il tempo caratteristico è equiparabile al periodo il condensatore ci si aspetta che non abbia il tempo di caricarsi
    \item Mentre se il periodo risulta notevolmente superiore il condensatore non dovrebbe essere in grado di seguire le oscillazioni
\end{itemize}
\subsection{\textbf{Misure acquisite}}
Siamo andati a verificare, acquisendo attraverso l'oscilloscopio i dati, come i grafici voltaggio tempo variassero
 a seconda della relazione tra il tempo caratteristico e il periodo dell'onda quadra. In seguito i grafici:

 [GRAFICI (mettere nella descrizione la condizione)]

Per il primo caso è possibile avere una stima rapida del tempo caratteristico, attraverso una legge che
 collega due punti consegutivi con la formula:

\[\tau= \frac{t_2-t_1}{\log{\frac{V_0-V_1}{V_0-V_2}}}\]

di seguito i dati acquisiti e il fit:

[TABELLA E FIT]

Per il secondo caso invece andiamo a verificare il seguente accordo teorico
, sia per l'onda quadra, sia per un'onda triangolare:

\[V_{output}=\frac{1}{\tau} \int V_{imput}(t) \, dt
\]

[CAPIRE COSA METTERE]

Infine si è deciso di analizzare il comportamento della tensione anche ai capi della reisistenza.
Pure in questo caso abbiamo analizzato il comportamento al variare della relazione tra il tempo caratteristico e il periodo dell'onda.

 [GRAFICI (mettere nella descrizione la condizione)]

Questi grafici sono quello che ci saremmo aspettati dalle considerazioni ideali.
\section{Circuito RLC}
\subsection{Introduzione}
In questo circuito dalla risoluzione dell'equazione differenziale abbiamo due possibili funzioni che esprimono la tensione
agli estremi della capacità, a seconda del fattore di Merito. Se: 
\[Q=\frac{\omega_0}{\Gamma}= \frac{1}{R}\sqrt{\frac{L}{C}}>\frac{1}{2}
\quad \Rightarrow \quad V_c(t)=V_0[1-e^{-\frac{\Gamma}{2}t}(\cos{(\omega^*t)}+\frac{\Gamma}{2\omega^*}\sin{(\omega^*t)})] \]
Ovvero siamo nel caso in cui si comporta come l'oscillatore sottosmorzato, dove all'aumentare del fattore di merito le oscillazioni risultano più ampie con uno smorzamento pià lento.

\subsection{\textbf{Misure acquisite}}
Nel nostro caso il fattore di merito ideale risulta:
\[Q=5 [VALOre A CASO] \]

Di conseguenza andiamo a verificare se il comportamenteo del circuito è come aspettato. 

[GRAFICOOOOO] 

Adesso da questo Grafico è possibile ricavare le seguenti stime effettive

\subsection{\textbf{Analisi dati}}

%\renewcommand{\thesection}{}
%\renewcommand{\thesubsection}{\arabic{subsection})}
\section{Conclusioni} %Conclusioni
[DA FARE PER FORZA DOPO LAB]
\end{document}